\emph{Pour $t \in [0~;~+\infty[$, $f(t) = 10 \times 1,15^{t}$ où $t$ est exprimé en heures et
$f(t)$ en milliers de bactéries.}

\begin{enumerate}
	\item
	\begin{enumerate}
		\item Au début de l'expérience, $t=0$.
		$f(0) = 10 \times 1,15^{0} = 10 \times 1 = 10$.

		Au début de l'expérience, la population compte 10 milliers de bactéries, soit
		\np{10000} bactéries.

		\item $f(1) = 10 \times 1,15 = 11,5$.

		Au bout d'une heure, la population compte 11,5 milliers de bactéries, soit
		\np{11500} bactéries.
	\end{enumerate}

	\item Une fonction de la forme $t \longmapsto k \times a^t$ avec $k > 0$ et $a > 0$ est croissante
	si $a > 1$ et décroissante si $0 < a < 1$.

	Ici $k = 10 > 0$ et $a = 1,15 > 1$ : la fonction $f$ est \textbf{croissante} sur $[0 ; +\infty[$.

	\item $t = 1,5$~h correspond à 1~h~30~min et $f(t)$ s'exprime en milliers de bactéries.

	Au bout d'une heure et demie, la population compte environ \np{12332} bactéries.
\end{enumerate}
